\section{Architecture Patterns}
We separate our system into three main layers:
\begin{itemize}
    \item \textbf{Presentation Layer:} Handles HTTP requests using controllers
    \item \textbf{Business Logic:} Contains core logic and services
    \item \textbf{Data Layer:} Manages database access
\end{itemize}
This separation keeps things clean and makes testing easier.
% To go further, we apply ideas from Clean Architecture:
% \begin{itemize}
% \item We keep business logic separate from frameworks
% \item We use ports and adapters to connect to databases and other systems
% \item Our use cases drive the app, not the UI or storage
% \end{itemize}
\subsection{Backend Patterns}
For working with data, we separate database details from domain logic using the Repository pattern. This makes our business logic simpler and more testable.

Communication happens through REST APIs where we separate resources clearly:
\begin{itemize}
    \item Endpoints match resources (e.g., \texttt{/users}, \texttt{/games})
    \item We follow HTTP standards and use JSON
    \item The API is stateless — each request stands on its own
\end{itemize}
For background tasks, we separate immediate actions from delayed ones using events. For example, when a score is saved, we trigger a separate event to update the leaderboard. This keeps different parts of the system independent.

We also separate commands (write operations) from queries (read operations). This CQRS pattern lets us optimize reads and writes differently.
\subsection{Frontend Patterns}
We separate the frontend into components. Each piece of UI is self-contained and can be reused.
% State is managed using a single store (Redux-style), so data flows in one direction: action → state → UI. This makes debugging easier and keeps logic consistent.
\subsection{System-Wide Practices}
We use dependency injection throughout the backend to separate what components do from how they're created. This lets us swap components easily and write better tests.
% When calling external services, we plan to use a circuit breaker pattern. It
% helps avoid system-wide failures when external dependencies are down.
\subsection{Deployment and Design Choices}
Right now, we use a modular monolith — all parts live in one deployable unit, but we separate the code into modules. This makes development simple, but we can split into microservices later if needed.
% We don't have an API gateway yet, but the structure allows us to add one in the future to manage routing, authentication, and frontend-specific APIs.
\subsection{Why This Works for Us}
By separating concerns this way, we get an architecture that is:
\begin{itemize}
    \item \textbf{Simple to understand}
    \item \textbf{Easy to test and maintain}
    \item \textbf{Ready to scale and evolve}
\end{itemize}
We focused on keeping the system flexible and clean by separating responsibilities without over-engineering at this stage.